\chapter{Conclusioni e sviluppi futuri}
\label{chap:conclusioni}

\section{Riepilogo dei contributi}

La replica dei primi due stadi della pipeline BoNet~\cite{escobar2019bonet},
hand detection e keypoint estimation, ha raggiunto prestazioni comparabili
a quelle del paper di riferimento (Capitolo~\ref{chap:risultati}), pur
addestrando i modelli da zero e su un sottoinsieme ridotto del dataset
RHPE, confermando la fattibilità dell'approccio anche in condizioni di
risorse limitate.

Il secondo obiettivo, la predizione dei keypoint su entrambe le mani, ha
portato al confronto di quattro varianti di modello e quattro strategie
di combinazione (Sez.~\ref{sec:esplorazione_dx_metodo}), da cui è emerso
\texttt{mirrored\_full} come soluzione più robusta: l'unico modello in
grado di produrre detection e keypoint per entrambe le mani in un unico
passaggio, senza necessitare di una bounding box esterna
(Sez.~\ref{sec:esplorazione_dx}, Capitolo~\ref{chap:discussione}).

\section{Sviluppi futuri}
% TODO: pseudo-labeling sulle 50k immagini ospedaliere, classificatore
% dx/sx, completamento del Bone Age Assessment (prendere spunto da SIMBA,
% Gonzalez et al. MICCAI 2020: riusa heatmap da keypoint + identity markers
% eta cronologica/genere via moltiplicatori appresi + relative bone age),
% training su dataset completo
